% ECONOMICS_PROBLEM: {"id": "D72-535", "title": "Economic origins and consequences of populism", "jel_code": "D72", "source_id": "OP-0124"}
% !TeX program = xelatex
\documentclass[11pt,a4paper]{article}
\usepackage[margin=23mm]{geometry}
\usepackage{fontspec}
\setmainfont{Latin Modern Roman}
\usepackage{amsmath,amssymb,mathtools}
\usepackage{xcolor,enumitem,booktabs,longtable,array,xurl,hyperref}
\definecolor{muted}{HTML}{53636C}
\hypersetup{colorlinks=true,urlcolor=blue,linkcolor=blue}
\setlength{\parindent}{0pt}
\setlength{\parskip}{5pt}
\newcommand{\R}{\mathbb R}
\newcommand{\E}{\mathbb E}
\newcommand{\N}{\mathbb N}
\newcommand{\ind}{\mathbf 1}
\newcommand{\Var}{\operatorname{Var}}
\newcommand{\Cov}{\operatorname{Cov}}
\newcommand{\supp}{\operatorname{supp}}
\DeclareMathOperator*{\argmin}{arg\,min}
\DeclareMathOperator*{\argmax}{arg\,max}
\newcommand{\entrymeta}[3]{{\sffamily\small\color{muted}\textbf{\texttt{#1}}\quad|\quad #2\par #3\par}}
\newcommand{\Problemref}[1]{\texttt{#1}}
\newcommand{\JELref}[2]{\texttt{#2} (JEL \texttt{#1})}
\begin{document}
\section*{Economic origins and consequences of populism}
\textbf{Problem ID:} \texttt{D72-535}\quad \textbf{JEL:} \texttt{D72}\par
% BEGIN ECONOMICS STATEMENT
\label{problem:OP-0124}
\label{atlas:prob:Q4}
{\sffamily\small\color{muted}\textbf{Q4} | Political economy | Editorial research problem family\par}
\subsection*{Definitions and assumptions}
Fix regions $r$, a horizon $h$, a preregistered definition of a populist party, and vote share $P_{rh}\in[0,1]$. Let the vector $a\in\mathcal A$ describe feasible trade, technology, migration, redistribution and identity-salience interventions. Let $P_{rh}(a)$ and $Y_{rh}(a)$ be potential vote shares and economic outcomes, allowing regional spillovers and strategic party entry. Fix a baseline $a^0$ and comparison $a^1$. For $K$ intervention components and $S\subseteq\{1,\ldots,K\}$, let $a^S$ adopt the comparison component for indices in $S$ and the baseline component otherwise; require these mixed regimes to be feasible.
\subsection*{Formal research question}
Identify or bound $v(S)=\mathbb E[P_{rh}(a^S)]$ and the subsequent economic contrast $\mathbb E[Y_{rh}(a^1)-Y_{rh}(a^0)]$. If a component attribution is required, specify the convention, for example the Shapley attribution
\[
\phi_k=\sum_{S\subseteq\{1,\ldots,K\}\setminus\{k\}}
\frac{|S|!(K-|S|-1)!}{K!}\,[v(S\cup\{k\})-v(S)].
\]
Under this definition $\sum_k\phi_k=v(\{1,\ldots,K\})-v(\varnothing)$. Determine which comparisons are identified under the declared model class, assignment mechanism and transport assumptions.
\subsection*{Interpretation and scope}
The decomposition is an editorial accounting convention for specified interventions, not a unique fraction of history caused by each mechanism. It accommodates interaction effects by averaging their allocation across orderings. If mixed interventions are infeasible, the Shapley target must be abandoned or replaced. Trade shocks do not manipulate culture independently, and a populist vote need not imply a particular economic policy. Electoral and policy consequences therefore remain distinct estimands.
\subsection*{Source audit and status}
Rodrik (2018) is verified, while the atlas link is the 2017 working paper. Colantone--Stanig (2018) is ambiguous between two verified papers; both are supplied. The undated phrase ``Autor et al. political China-shock research'' is resolved to a selected 2020 article, with its 2016 working-paper antecedent distinguished. These sources support mechanisms and specific evidence, not the editorial attribution formula or a universal causal decomposition.
\subsection*{References}
{\footnotesize\raggedright
\begin{enumerate}[label={\textbf{[S\arabic*]}},leftmargin=3em]
\item Dani Rodrik (2018). \emph{Populism and the Economics of Globalization}. Journal of International Business Policy 1, 12--33.\par
\textit{Locator and source role:} Abstract and globalization mechanisms. The linked NBER Working Paper 23559 is dated 2017, not 2018.\par
\url{https://www.hks.harvard.edu/publications/populism-and-economics-globalization-0}
\item Italo Colantone and Piero Stanig (2018). \emph{Global Competition and Brexit}. American Political Science Review 112(2), 201--218.\par
\textit{Locator and source role:} Abstract and trade-exposure design; one of two plausible works matching the original author-year citation.\par
\url{https://www.cambridge.org/core/journals/american-political-science-review/article/abs/global-competition-and-brexit/C843990101DB9232B654E77130F88398}
\item Italo Colantone and Piero Stanig (2018). \emph{The Trade Origins of Economic Nationalism: Import Competition and Voting Behavior in Western Europe}. American Journal of Political Science 62(4), 936--953.\par
\textit{Locator and source role:} Abstract; second plausible author-year match, retained to expose the original ambiguity.\par
\url{https://onlinelibrary.wiley.com/doi/10.1111/ajps.12358}
\item David Autor, David Dorn, Gordon Hanson, and Kaveh Majlesi (2020). \emph{Importing Political Polarization? The Electoral Consequences of Rising Trade Exposure}. American Economic Review 110(10), 3139--3183.\par
\textit{Locator and source role:} Abstract and local import-exposure identification; resolves the atlas phrase political China-shock research to a selected identifiable paper.\par
\url{https://www.aeaweb.org/articles?id=10.1257/aer.20170011}
\end{enumerate}
}

\subsection*{Common conventions: Source mathematical conventions}

Unless an entry states otherwise, agent and firm sets are finite. State and action spaces are measurable, strategies and outcome maps are measurable, probability laws are specified, and expectations exist for the targets under discussion. Infinite-horizon discount factors lie in $(0,1)$ and discounted objectives are finite. Norms, tolerances and complexity input sizes must be specified for computational comparisons. Constants or primitives that appear locally are inputs, not empirical facts asserted by this catalogue.

\textbf{Identification.} A statistical model $m\in\mathcal M$ implies an observable law $P_m$ and target $\theta(m)$. At law $P$, its identified set is
\[
\Theta(P;\mathcal M)=\{\theta(m):m\in\mathcal M,\ P_m=P\}.
\]
Point identification holds when this set is a singleton, or equivalently when $P_{m_1}=P_{m_2}$ implies $\theta(m_1)=\theta(m_2)$ for all compatible models. Sharp bounds mean bounds attained, or approached when the boundary is not attained, by compatible models. Writing this set is a definition, not a solution: the substantive task is to establish informative restrictions and derive or estimate the set. An unrestricted-confounding model may give only trivial bounds.

\textbf{Inference.} Identification, estimability and valid uncertainty quantification are separate questions. Fix $\alpha\in(0,1)$. A confidence procedure $C_n$ over a declared law class $\mathcal P$ has asymptotic uniform coverage at level $1-\alpha$ when
\[
\liminf_{n\to\infty}\inf_{P\in\mathcal P}\Pr_P\{\theta(P)\in C_n\}\geq1-\alpha.
\]
For set-identified targets the coverage event must be defined separately, for example coverage of the full identified set. If an entry uses identification-indexed models instead of uniquely indexed laws, the same coverage convention is applied over those models. Pointwise limits do not establish uniform validity, and asymptotic validity does not guarantee finite-sample precision. Sample sizes, dependence structures and weak-identification sequences are part of each application.

\textbf{Counterfactuals.} $Y_i(a)$ is a potential outcome under a specified intervention, including its timing, implementation and versions. It is not $Y_i$ conditional on voluntarily choosing $a$. A full intervention vector permits interference; shorthand scalar treatment notation does not silently impose no interference. Assignment, selection, attrition, anticipation and measurement must be specified. A claim about an instrument's local effect is not automatically a population policy effect.

\textbf{Economic equilibrium.} A counterfactual model includes best responses, constraints, feasibility, market clearing, state transitions and expectations consistent with the stated information. When several equilibria exist, either a declared selector is an input or the target is set-valued. Comparisons across policy regimes must use the stated selection convention. Reduced-form relationships are not presumed invariant after an equilibrium-changing intervention.

\textbf{Optimization and welfare.} Optimization is over the stated feasible set. If attainment is not established, $\sup$ or $\inf$ and $\varepsilon$-optimal choices replace an asserted maximizer or minimizer. For example, with value $V=\sup_{a\in\mathcal A}W(a)<\infty$, an $\varepsilon$-optimal set is $\{a\in\mathcal A:W(a)\geq V-\varepsilon\}$ for $\varepsilon>0$. Benefits, costs, risk and health or environmental outcomes can be aggregated only using declared units and conversion weights. Interpersonal utility weights, the treatment of fiscal transfers and foreign profits, discounting, and the behavioral welfare criterion are explicit inputs. A comparative-static derivative additionally requires existence and appropriate differentiability of the selected counterfactual.

\textbf{Sources.} Local labels [S1], [S2], and so forth restart in every question. Locators identify the relevant abstract, model, section or result at the level actually checked; a bibliographic or abstract-level verification is not represented as inspection of an entire proof. Working papers are distinguished from later publications and changing versions. Source summaries are paraphrased. All URL checks and editorial formulations refer to the audit date stated above.
% END ECONOMICS STATEMENT
\end{document}
